\section*{Hoofdstuk \ref{ch:g10:muscles-and-joints} --- Spieren en gewrichten}

\begin{solution}{exo:g10:muscles-and-joints:1}
\omterm{def:g10:muscles-and-joints:joint}{Kraakbeen} op de boteinden (soepel glijden, belasting verdelen); het
kapsel (sluit het \omterm{def:g10:muscles-and-joints:joint}{gewricht} af) met zijn gewrichtssmeer (smering);
\omterm{def:g10:muscles-and-joints:joint}{gewrichtsbanden} (houden de botten bijeen, beperken de bewegingen); in de
knie de menisci (verdelen de belasting).
\end{solution}

\begin{solution}{exo:g10:muscles-and-joints:2}
Een \omterm{def:g10:muscles-and-joints:muscle}{pees} verbindt een spier met een bot en brengt haar trek over; een
\omterm{def:g10:muscles-and-joints:joint}{gewrichtsband} verbindt bot met bot over een \omterm{def:g10:muscles-and-joints:joint}{gewricht} heen en beperkt de
beweging ervan. Geen van beide trekt samen.
\end{solution}

\begin{solution}{exo:g10:muscles-and-joints:3}
Een spier kan alleen trekken door korter te worden; duwen kan zij niet.
Om een \omterm{def:g10:muscles-and-joints:joint}{gewricht} beide kanten op te bewegen zijn een buiger en een
strekker nodig, die elk ongedaan maken wat de ander deed.
\end{solution}

\begin{solution}{exo:g10:muscles-and-joints:4}
Verzwikking: een \omterm{def:g10:muscles-and-joints:joint}{gewrichtsband} uitgerekt of gescheurd. Spierscheur:
\omterm{def:g10:muscles-and-joints:muscle}{spiervezels} gescheurd. Ontwrichting: een booteinde uit zijn \omterm{def:g10:muscles-and-joints:joint}{gewricht}
gedwongen, waarbij het kapsel en de \omterm{def:g10:muscles-and-joints:joint}{gewrichtsbanden} worden uitgerekt of
scheuren.
\end{solution}

\begin{solution}{exo:g10:muscles-and-joints:5}
Trage vezels: langzaam, zwak, onvermoeibaar, rijk aan \omterm{def:g10:cells-common-unit:organelle}{mitochondriën} en
haarvaten, rood, draaien op de ademhaling. Snelle vezels: snel, krachtig,
binnen een minuut moe, weinig \omterm{def:g10:cells-common-unit:organelle}{mitochondriën}, bleek, draaien op de
fosfaatvoorraad en de \omterm{def:g10:cell-metabolism:fermentation}{gisting}.
\end{solution}

\begin{solution}{exo:g10:muscles-and-joints:6}
Last $2 \times 9.8 = \qty{19.6}{N}$; kracht in de biceps $19.6 \times 32/4
\approx \qty{157}{N}$.
\end{solution}

\begin{solution}{exo:g10:muscles-and-joints:7}
Een duw kan uit beide richtingen komen, en elke spier kan er maar één
weerstaan. Trekken beide samen, dan verstijven zij het \omterm{def:g10:muscles-and-joints:joint}{gewricht} in beide
richtingen tegelijk; uit welke richting de duw ook komt, een van de twee
trekt er al tegenin.
\end{solution}

\begin{solution}{exo:g10:muscles-and-joints:8}
De vezel wordt \qty{0.6}{cm} korter; de \omterm{def:g10:muscles-and-joints:muscle}{pees} verplaatst het bot
\qty{0.6}{cm} op \qty{5}{cm} van de elleboog, dus verplaatst de hand zich
$0.6 \times 35/5 = \qty{4.2}{cm}$.
\end{solution}

\begin{solution}{exo:g10:muscles-and-joints:9}
Beweeglijkheid vergt een ondiepe kom en losse \omterm{def:g10:muscles-and-joints:joint}{gewrichtsbanden}; diezelfde
losheid laat de kop van het opperarmbeen de kom verlaten wanneer een
kracht hem voorbij het bereik duwt dat de banden kunnen houden.
\end{solution}

\begin{solution}{exo:g10:muscles-and-joints:10}
Uitdroging en zoutverlies door zweten, en vermoeidheid van de vezels aan
het eind van de wedstrijd. Remedie: de spier voorzichtig rekken en water
met zout drinken.
\end{solution}

\begin{solution}{exo:g10:muscles-and-joints:11}
Het dijbeen van de sprinter bestaat voor 70\% uit snelle vezels en dat
aandeel is grotendeels erfelijk; training vergroot vezels en verbetert
hun uitrusting, maar zet hun aantal niet om. De ongetrainde persoon
begint half om half en heeft meer trage vezels om te ontwikkelen.
\end{solution}

\begin{solution}{exo:g10:muscles-and-joints:12}
$4 \times 70 \times 9.8 \approx \qty{2.7}{kN}$. Spanning $2744/100
\approx \qty{27}{N/mm^2}$: ongeveer een kwart van de breukspanning, een
veiligheidsfactor vier.
\end{solution}

\begin{solution}{exo:g10:muscles-and-joints:13}
Spier is rijk van bloed voorzien en bevat reservecellen die vezels
opnieuw opbouwen: weken. \omterm{def:g10:muscles-and-joints:joint}{Gewrichtsbanden} hebben weinig vaten, dus komen
voedingsstoffen en herstelcellen er traag aan: maanden. \omterm{def:g10:muscles-and-joints:joint}{Kraakbeen} heeft
helemaal geen bloedvaten en vrijwel geen delende \omterm{def:g10:cells-common-unit:cell}{cellen}; het wordt alleen
door de gewrichtssmeer gevoed en herstelt nauwelijks.
\end{solution}

\begin{solution}{exo:g10:muscles-and-joints:14}
Bodybuilding verdikt de vezels (meer myofibrillen in elk) en rekken
verlengt de eenheid van spier en \omterm{def:g10:muscles-and-joints:muscle}{pees} en de tolerantie van het kapsel.
Het aantal vezels ligt vroeg in het leven vast; volwassen spieren groeien
of krimpen door de omvang van bestaande vezels te veranderen.
\end{solution}

\begin{solution}{exo:g10:muscles-and-joints:15}
Op de trap draagt het \omterm{def:g10:muscles-and-joints:joint}{gewricht} enkele keren het lichaamsgewicht; dunner
\omterm{def:g10:muscles-and-joints:joint}{kraakbeen} brengt die druk over op het bot eronder, dat zenuwen heeft en
pijn doet. Met minder \omterm{def:g10:muscles-and-joints:joint}{kraakbeen} zijn de menisci het laatste wat de
belasting verdeelt. Sterke spieren rond de knie vangen een deel van de
schok op en houden het \omterm{def:g10:muscles-and-joints:joint}{gewricht} stabiel, wat de piekdruk op het
resterende \omterm{def:g10:muscles-and-joints:joint}{kraakbeen} verlaagt.
\end{solution}

\begin{solution}{pb:g10:muscles-and-joints:1}
\textbf{1.} $70 \times 9.8 = \qty{686}{N}$.

\textbf{2.} $686 \times 10/5 \approx \qty{1.4}{kN}$.

\textbf{3.} $686 \times 20/5 \approx \qty{2.7}{kN}$.

\textbf{4.} Twee en vier keer het lichaamsgewicht: op elke trede trekt de
quadriceps met het gewicht van vier mensen, en daarom branden de dijen
bij een lange klim.

\textbf{5.} De hamstrings, aan de achterkant van het dijbeen; zij
ontspannen en worden uitgerekt terwijl de quadriceps de knie strekt, en
trekken licht samen om het \omterm{def:g10:muscles-and-joints:joint}{gewricht} stabiel te houden.

\textbf{6.} Grondkracht $3 \times 686 \approx \qty{2.06}{kN}$;
peeskracht $2058 \times 20/5 \approx \qty{8.2}{kN}$, twaalf keer het
lichaamsgewicht.

\textbf{7.} $8232/120 \approx \qty{69}{N/mm^2}$: ongeveer twee derde van
de breukspanning. Elke landing gebruikt het grootste deel van de marge
van de \omterm{def:g10:muscles-and-joints:muscle}{pees}.

\textbf{8.} Belasting $8.2 + 2.1 \approx \qty{10.3}{kN}$. Met de menisci:
$10300/600 \approx \qty{17}{N/mm^2}$; zonder: $10300/100 \approx
\qty{103}{N/mm^2}$, zes keer meer.

\textbf{9.} Zes keer de druk op hetzelfde \omterm{def:g10:muscles-and-joints:joint}{kraakbeen} bij elke landing: het
oppervlak wordt sneller vermorzeld dan het kan worden onderhouden, en het
slijt.

\textbf{10.} Een laagje vloeistof tussen de kraakbeenvlakken, zodat zij
tijdens het buigen zonder droge wrijving langs elkaar glijden, en een
lichte demping van de klap.

\textbf{11.} Een draai van het scheenbeen onder het dijbeen, met een
zijwaartse buiging: draaiingen waarvoor de knie niet is gebouwd en die de
\omterm{def:g10:muscles-and-joints:joint}{gewrichtsbanden} moeten weerstaan.

\textbf{12.} De quadriceps trekt alleen in de lengterichting van het
been en strekt de knie; tegen een draai kan zij niets. En de blessure
gebeurt in enkele honderdsten van een seconde, sneller dan enige
spierreactie.

\textbf{13.} Bloed uit de vaten van de gescheurde \omterm{def:g10:muscles-and-joints:joint}{gewrichtsband} en het
kapsel, plus extra gewrichtssmeer: het kapsel was doorbroken of ontstoken
en het \omterm{def:g10:muscles-and-joints:joint}{gewricht} liep vol.

\textbf{14.} De \omterm{def:g10:muscles-and-joints:joint}{gewrichtsband}, slecht van bloed voorzien, heeft maanden
nodig en vaak een operatie; de gescheurde meniscus, \omterm{def:g10:muscles-and-joints:joint}{kraakbeen} vrijwel
zonder vaten, geneest meestal niet en het gescheurde deel wordt
bijgeknipt of gehecht.

\textbf{15.} De hamstrings zijn de antagonisten van de quadriceps en
trekken het scheenbeen naar achteren; een sterk paar dat samen aanspant,
verstijft de knie tegen het naar voren glijden dat de gescheurde band
niet meer tegenhoudt, en beschermt het \omterm{def:g10:muscles-and-joints:joint}{gewricht} voor en na de operatie.

\textbf{16.} Ongeveer 4000 keer per dag bij \qty{1.4}{kN} of meer.

\textbf{17.} Ongeveer $\num{1.5e6}$ belastingen per jaar. Een
peesontsteking is schade die zich sneller ophoopt dan het trage herstel
van peesweefsel; alleen rust laat het herstel bijkomen, want kracht
verlaagt de belasting van elke stap niet.

\textbf{18.} Landingen $40 \times 8.2 \approx \qty{330}{kN}$ in totaal;
stappen $4000 \times 1.4 \approx \qty{5600}{kN}$ --- de stappen dragen
zeventien keer meer totale belasting, maar de landingen dragen de pieken
die met scheuren dreigen.

\textbf{19.} De grondkracht verdubbelen verdubbelt bij elke pas de
krachten in de \omterm{def:g10:muscles-and-joints:muscle}{pees} en het \omterm{def:g10:muscles-and-joints:joint}{gewricht}; fietsen en zwemmen laten de spieren
werken terwijl zij de knie ver onder het lichaamsgewicht belasten.

\textbf{20.} De \omterm{def:g10:muscles-and-joints:muscle}{pees} van de quadriceps draagt vier keer het
lichaamsgewicht op een trap en ongeveer twaalf keer bij een landing; de
\omterm{def:g10:muscles-and-joints:joint}{gewrichtsbanden}, die door die belastingen nooit worden beproefd, waren het
weefsel dat de draai vond.
\end{solution}
